%! Piecewise-Defined Functions — Unit 2, Lesson 2.6 — Group Activity (Answer Key)
\documentclass[10pt]{article}
\usepackage{atda-article}
\usepackage{atda-key}   % pulls in -boxes; do NOT also load -boxes
\newcommand{\TallMath}[1]{$\displaystyle #1\rule[-1.4em]{0pt}{3.2em}$}
\setlength{\workrowsep}{6pt}

\begin{document}

\pageheader{Unit 2, Lesson 2.6}{Group Activity --- Answer Key}
% NAMESTRIP: no \namedateperiod here — mirrors the blank. Teacher-only prose goes
% in the lesson plan as \begin{teachernote}[Group Activity], never in a key.

\vspace{0.15in}

\begin{scenariobox}[The Scenario]{royal}
\small
Every question today is answered by reading a graph. \textbf{The rule for the whole activity: before
you use a rule or read a value, say out loud which piece owns the input.} And whenever a boundary
point is involved, \textbf{let the dots decide} --- filled means included, open means not.
\textbf{Everyone starts at Tier R.} When your whole group can do Tier R without help, move on.
\end{scenariobox}

\vspace{0.10in}

\boxguard[30]
\begin{tcolorbox}[colback=white, colframe=black!40, title=\textbf{Tier R --- Remediate},
    fonttitle=\bfseries, arc=2mm, left=3mm, right=3mm, top=2mm, bottom=2mm, breakable]
\small
A phone plan costs \$30 a month for anything up to $10$ gigabytes of data, and \$5 for each gigabyte
beyond $10$. The graph shows $C(g)$, the monthly cost in dollars when $g$ gigabytes are used. The plan
shuts data off at $20$ gigabytes.

\begin{center}
\begin{tikzpicture}
\begin{axis}[
    width=10.4cm, height=4.6cm,
    xlabel={$g$ (gigabytes used)}, ylabel={$C$ (monthly cost, \$)},
    xmin=0, xmax=20, ymin=0, ymax=90,
    xtick={0,2,4,6,8,10,12,14,16,18,20}, ytick={0,10,20,30,40,50,60,70,80},
    minor x tick num=1, minor y tick num=1,
    grid=both, grid style={linegray!45}, minor grid style={linegray!30},
    axis lines=left,
    tick label style={font=\tiny}, label style={font=\scriptsize}]
  \addplot[royal, very thick] coordinates {(0,30) (10,30) (20,80)};
  \addplot[keyred, only marks, mark=*, mark size=1.7pt]
    coordinates {(4,30) (10,30) (14,50) (18,70) (20,80)};
\end{axis}
\end{tikzpicture}
\end{center}

\begin{enumerate}[leftmargin=1.6em, itemsep=4pt, topsep=2pt]
  \item Read the graph and fill in both rows. For the second row write \emph{flat} or \emph{climbing}
        to say which piece owns that input.
\smallskip

\renewcommand{\arraystretch}{1.3}
\begin{tabular}{@{}|l|c|c|c|c|c|@{}}
\hline
$g$ (gigabytes) & $4$ & $10$ & $14$ & $18$ & $20$ \\ \hline
$C(g)$ (\$) & \ans{$30$} & \ans{$30$} & \ans{$50$} & \ans{$70$} & \ans{$80$} \\ \hline
Which piece? & \ans{flat} & \ans{flat} & \ans{climb.} & \ans{climb.} & \ans{climb.} \\
\hline
\end{tabular}

  \item Give the \emph{boundary point} --- the input where one rule hands off to the next --- and say
        in one sentence what it means to a customer.

\ansline{The boundary point is $g=10$ gigabytes. Up to and including $10$ GB the bill is a flat \$30 no}\\
\ansline{matter how little you use; past $10$ GB every extra gigabyte adds \$5 to that same \$30 base.}\\

  \item Name the interval of inputs on which $C$ is constant and the interval on which $C$ is
        increasing. Then say whether the graph is continuous, and how you can tell from the picture.

\ansline{Constant on $[0,10]$; increasing on $[10,20]$.}\\
\ansline{Yes, continuous. Both pieces meet at the same point $(10,30)$ --- there is no gap and no}\\
\ansline{open circle, so the whole graph can be traced without lifting the pencil.}\\

  \item Give the domain and the range of $C$. Then give the absolute maximum and the absolute
        minimum, each with a \emph{value and a location}. Read the minimum's location carefully ---
        it is not a single point.

\ansline{Domain $[0,20]$ gigabytes; range $[30,80]$ dollars.}\\
\ansline{Absolute maximum \$80 at $g=20$ --- a single point, at the right endpoint.}\\
\ansline{Absolute minimum \$30, attained at \emph{every} $g$ from $0$ to $10$: the whole interval $[0,10]$.}\\
\end{enumerate}
\end{tcolorbox}

\vspace{0.10in}

\boxguard[30]
\begin{tcolorbox}[colback=white, colframe=black!40,
    title=\textbf{Tier A --- Approaching Proficiency},
    fonttitle=\bfseries, arc=2mm, left=3mm, right=3mm, top=2mm, bottom=2mm, breakable]
\small
Graph (a) shows $G(h)$, the cost in dollars of parking in a city garage for $h$ hours: \$4 for up to
and including $1$ hour, \$7 for over $1$ up to and including $3$, and \$12 for over $3$ up to and
including $6$. Graph (b) shows a function $f$ with no context attached. \textbf{In both graphs, filled
circles are part of the graph and open circles are not.}

\begin{center}
\begin{tabular}{@{}c@{\hspace{0.5cm}}c@{}}
\begin{tikzpicture}
\begin{axis}[
    width=6.5cm, height=4.4cm, title={\scriptsize (a) $G$, garage cost},
    xlabel={$h$ (hours parked)}, ylabel={$G$ (\$)},
    xmin=0, xmax=6.4, ymin=0, ymax=15,
    xtick={0,1,2,3,4,5,6}, ytick={0,2,4,6,8,10,12,14},
    grid=both, grid style={linegray!45},
    axis lines=left,
    tick label style={font=\tiny}, label style={font=\tiny}]
  \addplot[royal, very thick] coordinates {(0,4) (1,4)};
  \addplot[royal, very thick] coordinates {(1,7) (3,7)};
  \addplot[royal, very thick] coordinates {(3,12) (6,12)};
  \addplot[royal, only marks, mark=*, mark size=1.9pt] coordinates {(1,4) (3,7) (6,12)};
  \addplot[royal, only marks, mark=*, mark size=1.9pt, mark options={fill=white}]
    coordinates {(0,4) (1,7) (3,12)};
\end{axis}
\end{tikzpicture}
&
\begin{tikzpicture}
\begin{axis}[
    width=6.3cm, height=4.4cm, title={\scriptsize (b) $f$},
    xmin=-3, xmax=7, ymin=-3, ymax=7,
    xtick={-3,-2,-1,0,1,2,3,4,5,6,7}, ytick={-2,0,2,4,6},
    minor y tick num=1,
    grid=both, grid style={linegray!45}, minor grid style={linegray!30},
    axis lines=left, tick label style={font=\tiny}]
  \addplot[royal, very thick, domain=-2:2, samples=60]{x^2};
  \addplot[royal, very thick] coordinates {(2,6) (6,-2)};
  \addplot[royal, only marks, mark=*, mark size=1.9pt] coordinates {(-2,4) (2,4) (6,-2)};
  \addplot[royal, only marks, mark=*, mark size=1.9pt, mark options={fill=white}]
    coordinates {(2,6)};
\end{axis}
\end{tikzpicture}
\end{tabular}
\end{center}

\begin{enumerate}[leftmargin=1.6em, itemsep=4pt, topsep=2pt]
  \item \textbf{Graph (a).} Fill in the costs, then answer the boundary question underneath.
\smallskip

\renewcommand{\arraystretch}{1.3}
\begin{tabular}{@{}|l|c|c|c|c|c|@{}}
\hline
$h$ (hours) & $1$ & $2$ & $3$ & $4$ & $6$ \\ \hline
$G(h)$ (\$) & \ans{$4$} & \ans{$7$} & \ans{$7$} & \ans{$12$} & \ans{$12$} \\
\hline
\end{tabular}
\smallskip

A driver parks for \emph{exactly} $3$ hours and is charged \$12. Using the dots, say what they should
have been charged and which circle proves it.

\ansline{They should have been charged \$7. The \emph{filled} circle at $(3,7)$ is on the graph, so $G(3)=7$;}\\
\ansline{the circle at $(3,12)$ is \emph{open}, so $12$ is not the cost at $3$ hours --- only past $3$ hours.}\\

  \item \textbf{Graph (a).} Give the range of $G$. Then say whether $G$ is continuous, naming every
        input where you would have to lift your pencil.

\ansline{Range: the list $\{4, 7, 12\}$ --- only three costs are ever charged, so no interval describes it.}\\
\ansline{Discontinuous: you must lift the pencil at $h=1$ and at $h=3$. Two jumps, of \$3 and of \$5.}\\

  \item \textbf{Graph (a), the one that sounds wrong at first.} The cost clearly goes up as you park
        longer. Is there \emph{any} interval on which $G$ is increasing? Answer, and explain using
        what the graph does on each piece.

\ansline{No --- there is no interval on which $G$ is increasing. On each of its three pieces the graph}\\
\ansline{is \emph{constant}. The cost does rise, but it rises \emph{at} the two jumps, and a jump happens}\\
\ansline{at a single input, not across an interval. $G$ is constant on $(0,1]$, $(1,3]$ and $(3,6]$.}\\

  \item \textbf{Graph (b).} Give $f(-2)$, $f(0)$, $f(2)$, $f(4)$ and $f(6)$. Then give the domain of
        $f$, its zeros, and its $y$-intercept.

\ansline{$f(-2)=4$, \ $f(0)=0$, \ $f(2)=4$ (the \emph{filled} circle), \ $f(4)=2$, \ $f(6)=-2$.}\\
\ansline{Domain $[-2,6]$. Two zeros: $x=0$ and $x=5$. $y$-intercept $(0,0)$.}\\

  \item \textbf{Graph (b), the stretch.} Give the range of $f$, then give the absolute maximum and
        the absolute minimum. \emph{Look at the open circle before you answer} --- one of these two
        does not exist, and Lesson 2.5 gave you the words for why.

\ansline{Range $[-2,6)$ --- a parenthesis at $6$, because the circle at $(2,6)$ is open.}\\
\ansline{There is \textbf{no absolute maximum}: outputs just past $x=2$ get as close to $6$ as you like}\\
\ansline{and never reach it, so no output is the highest. Absolute minimum $-2$ at $x=6$.}\\
\end{enumerate}
\end{tcolorbox}

\vspace{0.10in}

\boxguard[30]
\begin{tcolorbox}[colback=white, colframe=black!40, title=\textbf{Tier E --- Extension},
    fonttitle=\bfseries, arc=2mm, left=3mm, right=3mm, top=2mm, bottom=2mm, breakable]
\small

\begin{enumerate}[leftmargin=1.6em, itemsep=4pt, topsep=2pt]
  \item \textbf{Two claims, and each one misreads a piece.} Name the mistake and give the correct
        statement.
        \begin{enumerate}[label=(\alph*), leftmargin=1.6em, itemsep=3pt, topsep=3pt]
          \item About the shipping graph in your notes: \emph{``At $w=2$ the cost is both \$6 and
                \$10, so $S$ is not a function.''}

\ansline{The open circle was counted as part of the graph. Above $w=2$ there is exactly one \emph{filled}}\\
\ansline{circle, at \$6, so $S(2)=\$6$ and $S$ is a function --- one output for that input, like every other.}\\

          \item About the garden centre's pay: \emph{``A $46$-hour week pays $46 \times \$12$ plus
                $6 \times \$18$, which is \$660.''}

\ansline{The first rule was applied to inputs it does not own. It covers only hours $1$--$40$; hours $41$--$46$}\\
\ansline{are paid at \$18 \emph{instead of} \$12. Correct: $40(12) + 6(18) = 480 + 108 = \$588$.}\\
        \end{enumerate}

  \item \textbf{Possible or impossible?} For each description, either name a graph from today that
        does it, or explain why nothing can.
        \begin{enumerate}[label=(\alph*), leftmargin=1.6em, itemsep=3pt, topsep=3pt]
          \item A piecewise function that is continuous everywhere and still has a sharp corner.

\ansline{Possible --- the paycheck $P$ at $h=40$, and the drone at $t=9$. A corner is not a break.}\\

          \item A piecewise function whose graph gives two different outputs at one input.

\ansline{Impossible. The pieces are not allowed to overlap, so each input belongs to exactly one piece.}\\
\ansline{Two outputs at one input is not a function at all --- which is what the open circles prevent.}\\

          \item A piecewise function that is constant on every one of its pieces, yet whose outputs
                are larger at the right end than at the left.

\ansline{Possible --- the shipping cost $S$ and the garage cost $G$ both do it. Every piece is flat, and}\\
\ansline{all of the rising happens \emph{at} the jumps between pieces rather than along any piece.}\\
        \end{enumerate}

  \item \textbf{Write the rules.} A pool charges a flat \$3 to enter and then \$2 for each hour after
        the second. The graph of $h(x)$, the total cost in dollars for $x$ hours, is below. Write the
        two rules \emph{and} the stretch of inputs each one owns, then check your second rule by
        computing the cost at $x=6$ and comparing it with the graph.

\begin{center}
\begin{tikzpicture}
\begin{axis}[
    width=8.4cm, height=3.6cm,
    xlabel={$x$ (hours)}, ylabel={$h$ (total cost, \$)},
    xmin=0, xmax=6, ymin=0, ymax=13,
    xtick={0,1,2,3,4,5,6}, ytick={0,3,5,7,9,11},
    grid=both, grid style={linegray!45},
    axis lines=left,
    tick label style={font=\tiny}, label style={font=\tiny}]
  \addplot[royal, very thick] coordinates {(0,3) (2,3) (6,11)};
  \addplot[keyred, only marks, mark=*, mark size=1.7pt]
    coordinates {(0,3) (2,3) (4,7) (6,11)};
\end{axis}
\end{tikzpicture}
\end{center}

\ansline{First piece: $h(x)=3$, owning $0 \le x \le 2$, that is $[0,2]$ --- the flat entry fee.}\\
\ansline{Second piece: $h(x)=3+2(x-2)$, owning $2 < x \le 6$, that is $(2,6]$ --- \$2 an hour after hour $2$.}\\
\ansline{Check at $x=6$: $3+2(6-2)=3+8=11$, and the graph reads \$11 at $x=6$. The rules match the picture.}\\
\ansline{Note the boundary: both rules give \$3 at $x=2$, so the graph has a corner there but no break.}\\
\end{enumerate}
\end{tcolorbox}

\end{document}
